\section{Limits of sequences}
\label{sec:limits-of-sequences}
A sequence assigns a real number $s_n$ to each index $n$. Unless the
exercise states otherwise, the indices here are $n=1,2,3,\ldots$.
The order of the terms matters, even when their values repeat.

We use Ross's Definition 7.1: $s_n\to s$ means that, for every
$\varepsilon>0$, there is a number $N$ such that
\[
 n>N\quad\Longrightarrow\quad \abs{s_n-s}<\varepsilon
 \qquad(n\in\N).
\]
The choice of $N$ may depend on $\varepsilon$ but cannot depend on
the later index $n$. A real threshold $N$ is permitted; the
Archimedean property ensures that there are integer indices past it.
Convergence here means convergence to a finite real number.
Divergence means that no such real limit exists.

Although several exercises in this section ask only for answers,
we include the algebra and estimates that explain those answers.
The proofs use elementary inequalities and the limit definition,
rather than the limit laws introduced in Section 9.

\exercise{7.1}{The first five terms of four sequences}
\goal Compute the terms with indices $n=1,2,3,4,5$ for
\[
 \text{(a)}\ s_n=\frac1{3n+1},\qquad
 \text{(b)}\ b_n=\frac{3n+1}{4n-1},\qquad
 \text{(c)}\ c_n=\frac n{3^n},\qquad
 \text{(d)}\ d_n=\sin\left(\frac{n\pi}{4}\right).
\]

\idea Substitute the index into every occurrence of $n$, evaluate
the numerator and denominator separately, and simplify where possible.
Indices start at $1$. In part \textup{(a)}, the denominator $3n+1$
means $3$ times $n$, followed by addition of $1$.

\step{1. Part (a): A reciprocal with a linear denominator}
For $n\in\N$, we have $3n+1\ge4>0$, so all the indicated fractions
are defined. The five substitutions give
\begin{align*}
 s_1&=\frac1{3\cdot1+1}=\frac1{3+1}=\frac14,\\
 s_2&=\frac1{3\cdot2+1}=\frac1{6+1}=\frac17,\\
 s_3&=\frac1{3\cdot3+1}=\frac1{9+1}=\frac1{10},\\
 s_4&=\frac1{3\cdot4+1}=\frac1{12+1}=\frac1{13},\\
 s_5&=\frac1{3\cdot5+1}=\frac1{15+1}=\frac1{16}.
\end{align*}

\step{2. Part (b): A quotient of two linear expressions}
Here $4n-1\ge3>0$ for every $n\ge1$, so there is again no zero
denominator. Compute both parts of each fraction:
\begin{align*}
 b_1&=\frac{3\cdot1+1}{4\cdot1-1}
      =\frac{3+1}{4-1}=\frac43,\\
 b_2&=\frac{3\cdot2+1}{4\cdot2-1}
      =\frac{6+1}{8-1}=\frac77=1,\\
 b_3&=\frac{3\cdot3+1}{4\cdot3-1}
      =\frac{9+1}{12-1}=\frac{10}{11},\\
 b_4&=\frac{3\cdot4+1}{4\cdot4-1}
      =\frac{12+1}{16-1}=\frac{13}{15},\\
 b_5&=\frac{3\cdot5+1}{4\cdot5-1}
      =\frac{15+1}{20-1}=\frac{16}{19}.
\end{align*}

\step{3. Part (c): A linear numerator and an exponential denominator}
The denominator is the positive integer power $3^n$. In particular,
$3^1=3$, $3^2=9$, $3^3=27$, $3^4=81$, and $3^5=243$.
Substitution yields
\begin{align*}
 c_1&=\frac1{3^1}=\frac13,\\
 c_2&=\frac2{3^2}=\frac29,\\
 c_3&=\frac3{3^3}=\frac3{27}=\frac19,\\
 c_4&=\frac4{3^4}=\frac4{81},\\
 c_5&=\frac5{3^5}=\frac5{243}.
\end{align*}
The cancellation in the third term divides numerator and denominator
by $3$; it does not change the value of that term.

\step{4. Part (d): Sine at multiples of an eighth-turn angle}
On the unit circle at angle $\pi/4$, the two coordinates are equal
and positive. If each is $t$, then $t^2+t^2=1$, so
$t^2=1/2$ and $t=\sqrt2/2$. Thus $\sin(\pi/4)=\sqrt2/2$.
The unit-circle values at $\pi/2$ and $\pi$ are $1$ and $0$.
Using these values and the identities
$\sin(\pi-\theta)=\sin\theta$ and
$\sin(\pi+\theta)=-\sin\theta$, we obtain
\begin{align*}
 d_1&=\sin\left(\frac{1\pi}{4}\right)
      =\sin\left(\frac\pi4\right)=\frac{\sqrt2}{2},\\
 d_2&=\sin\left(\frac{2\pi}{4}\right)
      =\sin\left(\frac\pi2\right)=1,\\
 d_3&=\sin\left(\frac{3\pi}{4}\right)
      =\sin\left(\pi-\frac\pi4\right)=\frac{\sqrt2}{2},\\
 d_4&=\sin\left(\frac{4\pi}{4}\right)=\sin\pi=0,\\
 d_5&=\sin\left(\frac{5\pi}{4}\right)
      =\sin\left(\pi+\frac\pi4\right)=-\frac{\sqrt2}{2}.
\end{align*}

\answer{The requested five-term lists are
\[
 \begin{array}{c|ccccc}
 \text{Part}&1&2&3&4&5\\ \hline
 \text{(a)}&\frac14&\frac17&\frac1{10}&\frac1{13}&\frac1{16}\\[3pt]
 \text{(b)}&\frac43&1&\frac{10}{11}&\frac{13}{15}&\frac{16}{19}\\[3pt]
 \text{(c)}&\frac13&\frac29&\frac19&\frac4{81}&\frac5{243}\\[3pt]
 \text{(d)}&\frac{\sqrt2}{2}&1&\frac{\sqrt2}{2}&0&-\frac{\sqrt2}{2}
 \end{array}
\]}

\exercise{7.2}{Convergence of the four preceding sequences}
\goal Decide which sequences in Exercise~7.1 converge and give
each finite limit when it exists.

\idea Use exact error estimates for the first three sequences.
For the sine sequence, find values $1$ and $-1$ at arbitrarily large
indices, then compare them with any proposed common limit.
The arguments use the limit definition directly: $u_n\to L$ means
that for every $\varepsilon>0$ there is a threshold $N$ such that
$n>N$ implies $\abs{u_n-L}<\varepsilon$.

\step{1. Part (a): The limit of $1/(3n+1)$}
For every positive integer $n$,
\[
 \abs{\frac1{3n+1}-0}
 =\frac1{3n+1}<\frac1{3n}.
\]
The reciprocal inequality follows because $3n+1>3n>0$.
Let $\varepsilon>0$. By the Archimedean property, choose an integer
$N>1/(3\varepsilon)$. If $n>N$, then $3n>1/\varepsilon$ and
\[
 \abs{\frac1{3n+1}-0}<\frac1{3n}<\varepsilon.
\]
This verifies the definition and proves $s_n\to0$.

\step{2. Part (b): The limit of $(3n+1)/(4n-1)$}
The ratio of the leading coefficients suggests $3/4$.
Compute the error over a common denominator to verify it:
\begin{align*}
 b_n-\frac34
 &=\frac{3n+1}{4n-1}-\frac34\\
 &=\frac{4(3n+1)-3(4n-1)}{4(4n-1)}\\
 &=\frac{12n+4-12n+3}{4(4n-1)}\\
 &=\frac7{4(4n-1)}.
\end{align*}
The denominator is positive for $n\ge1$, so this is also the
absolute error. Moreover,
$4n-1-3n=n-1\ge0$, hence $4n-1\ge3n$. Therefore
\[
 \abs{b_n-\frac34}
 =\frac7{4(4n-1)}\le\frac7{12n}.
\]
Given $\varepsilon>0$, choose an integer
$N>\max\{1,7/(12\varepsilon)\}$. For $n>N$ we get
\[
 \abs{b_n-\frac34}\le\frac7{12n}<\varepsilon.
\]
Thus $b_n\to3/4$. No quotient limit law is needed for this conclusion.

\step{3. Part (c): Bound the exponential denominator by a square}
We first prove $3^n\ge n^2$ for every integer $n\ge2$.
For $n=2$, this reads $9\ge4$. Suppose it holds for an integer
$n\ge2$. Then
\[
 3^{n+1}=3\cdot3^n\ge3n^2.
\]
Expanding the next square gives
\[
 3n^2-(n+1)^2
 =3n^2-(n^2+2n+1)
 =2n(n-1)-1\ge3>0,
\]
where $n\ge2$ gives $2n(n-1)\ge4$.
Hence $3n^2\ge(n+1)^2$, so $3^{n+1}\ge(n+1)^2$.
Mathematical induction proves the claimed bound.

Because all denominators are positive, the bound implies, for $n\ge2$,
\[
 0<\frac n{3^n}\le\frac n{n^2}=\frac1n.
\]
Let $\varepsilon>0$ and choose an integer
$N>\max\{2,1/\varepsilon\}$. Then $n>N$ gives
\[
 \abs{\frac n{3^n}-0}
 =\frac n{3^n}\le\frac1n<\varepsilon.
\]
It follows directly that $c_n\to0$.

\step{4. Part (d): Values separated by $2$ persist arbitrarily far out}
For every nonnegative integer $k$, periodicity of sine gives
\begin{align*}
 d_{8k+2}
 &=\sin\left(\frac{(8k+2)\pi}{4}\right)
  =\sin\left(2k\pi+\frac\pi2\right)=1,\\
 d_{8k+6}
 &=\sin\left(\frac{(8k+6)\pi}{4}\right)
  =\sin\left(2k\pi+\frac{3\pi}{2}\right)=-1.
\end{align*}
Both families of indices can be made larger than any given real
threshold $N$: choose an integer $k>\max\{0,N/8\}$.
Then $8k+2>N$ and $8k+6>N$.

Suppose, for a contradiction, that $d_n\to L\in\R$.
Use $\varepsilon=1/2$ in the definition. There is a threshold $N$
such that $n>N$ implies $\abs{d_n-L}<1/2$.
Choose $k$ so that both indices above exceed this $N$. Then
\[
 \abs{1-L}<\frac12,
 \qquad \abs{-1-L}<\frac12.
\]
The triangle inequality, together with
$\abs{L+1}=\abs{-1-L}$, would imply
\[
 2=\abs{1-(-1)}
   =\abs{(1-L)+(L+1)}
   \le\abs{1-L}+\abs{L+1}
   <\frac12+\frac12=1.
\]
This contradicts $2>1$. Therefore the sine sequence has no real limit.

\answer{Parts \textup{(a)}, \textup{(b)}, and \textup{(c)} converge
to $0$, $3/4$, and $0$, respectively. Part \textup{(d)} diverges.}

\exercise{7.3}{A catalogue of twenty sequences and their limits}

\goal Determine convergence or divergence for every sequence in
parts (a)--(t), and give the real limit whenever one exists.
Although the exercise does not require proofs, we include a justification
of every answer. The arguments use the definition of convergence
directly, rather than the limit laws developed later in Section~9.

\idea For a convergent sequence, subtract the proposed limit and
bound the resulting error by a reciprocal that becomes arbitrarily
small. For a divergent sequence, find either arbitrarily late values
that are arbitrarily large, or two arbitrarily late values separated
by a fixed positive distance.

\subsubsection*{Elementary facts used in the individual parts}

\step{Fact 1: Reciprocal error estimates imply convergence.}
The definition says that $u_n\to L$ if, for every $\varepsilon>0$,
there is a threshold $N$ such that
\[
 n>N\quad\Longrightarrow\quad\abs{u_n-L}<\varepsilon.
\]
Suppose we prove $\abs{u_n-L}\leq C/n$ for every $n\geq n_0$,
where $C>0$ is fixed. Given $\varepsilon>0$, choose a positive
integer $N>\max\{n_0,C/\varepsilon\}$, using the Archimedean
property. Then $n>N$ implies both $n\geq n_0$ and
$n>C/\varepsilon$. Multiplying the latter inequality by
$\varepsilon>0$ and dividing by $n>0$ gives
\[
 \abs{u_n-L}\leq\frac Cn<\varepsilon.
\]
This proves convergence from the definition. In particular,
$1/n\to0$. If an error is identically $0$, it is less than every
$\varepsilon>0$ without any additional estimate.

We will also use a shifted version. If, for all sufficiently large
$n$, $\abs{u_n-L}\leq C/(n-r)$, with fixed real $r$ and $n>r$,
choose $N$ larger than the starting index and $r+C/\varepsilon$.
Then $n>N$ gives $n-r>C/\varepsilon$, and the same calculation
gives $\abs{u_n-L}<\varepsilon$. This proves the shifted version
without any new limit theorem.

Only the tail is used in either proof. Thus a finite number of unusual
initial values does not affect the conclusion: increase $N$ to exceed
all their indices, and every $n>N$ still satisfies the required
estimate. This observation will matter when an initial denominator
is negative but never zero.

\step{Fact 2: Powers of two dominate a simple linear expression.}
For every integer $m\geq0$,
\[
 2^m\geq m+1.
\]
At $m=0$, both sides equal $1$. If the inequality holds at $m$,
then multiplication by $2>0$ gives
\[
 2^{m+1}=2\cdot2^m\geq2(m+1)=2m+2\geq m+2,
\]
where the last difference is $(2m+2)-(m+2)=m\geq0$.
Induction proves the assertion. Taking reciprocals of positive
quantities consequently gives
\[
 0<2^{-m}\leq\frac1{m+1}.
\]

\step{Fact 3: Separated arbitrarily late values rule out a real limit.}
Suppose some fixed $D>0$ has this property: for every threshold
$N$, there are indices $p,q>N$ with $\abs{u_p-u_q}\geq D$.
Assume, to obtain a contradiction, that $u_n\to L\in\R$.
Apply the definition with $\varepsilon=D/3$. It supplies $N$ such
that both late terms obey
$\abs{u_p-L}<D/3$ and $\abs{u_q-L}<D/3$. The triangle inequality
would then give
\[
 \abs{u_p-u_q}
 =\abs{(u_p-L)+(L-u_q)}
 \leq\abs{u_p-L}+\abs{L-u_q}
 <\frac D3+\frac D3=\frac{2D}{3}<D,
\]
contradicting the stipulated separation. Thus such a sequence
does not converge to a real number. This argument does not assume
any theorem about subsequences.

\step{Fact 4: Arbitrarily large late values also rule out a real limit.}
Suppose that for every $B\in\R$ and every threshold $N$, some
$n>N$ satisfies $u_n>B$. If $u_n\to L\in\R$, the definition
with $\varepsilon=1$ would supply $N$ for which
$\abs{u_n-L}<1$ whenever $n>N$. Choose a late value with
$u_n>L+1$. It satisfies $u_n-L>1$, and therefore
$\abs{u_n-L}=u_n-L>1$, a contradiction. This is a direct
nonconvergence argument; it does not invoke a later boundedness
theorem for convergent sequences.

\subsubsection*{Solutions of parts (a)--(t)}

\step{(a) The sequence $a_n=n/(n+1)$.}
For every $n\in\N$, $n+1\geq2>0$, so every term is defined.
To test the candidate $1$, subtract it using the common denominator
$n+1$:
\[
 a_n-1=\frac n{n+1}-\frac{n+1}{n+1}
 =\frac{n-(n+1)}{n+1}
 =-\frac1{n+1}.
\]
It follows that
\[
 \abs{a_n-1}=\frac1{n+1}\leq\frac1n.
\]
The inequality uses $n+1\geq n>0$ and the reversal of order for
positive reciprocals. Fact 1 with $C=1$ now proves the limit.

\answer{Part (a) converges, and $\displaystyle\lim_{n\to\infty}\frac n{n+1}=1$.}

\step{(b) The sequence $b_n=(n^2+3)/(n^2-3)$.}
The denominator is nonzero at every positive integer. At $n=1$
it is $1-3=-2$, while $n\geq2$ implies $n^2\geq4$ and
$n^2-3\geq1>0$. Thus there is no undefined term, including in
the initial part of the sequence.

Subtracting the candidate $1$ gives
\[
 b_n-1
 =\frac{n^2+3-(n^2-3)}{n^2-3}
 =\frac{n^2+3-n^2+3}{n^2-3}
 =\frac6{n^2-3}.
\]
For $n\geq3$, we have $n^2\geq9\geq6$, so $3\leq n^2/2$.
Subtracting $3$ from $n^2$ therefore gives
\[
 n^2-3\geq n^2-\frac{n^2}{2}=\frac{n^2}{2}>0.
\]
Consequently
\[
 \abs{b_n-1}
 =\frac6{n^2-3}
 \leq\frac6{n^2/2}
 =\frac{12}{n^2}
 \leq\frac{12}{n}.
\]
The last inequality follows from $n^2\geq n>0$ for $n\geq1$.
Fact 1 with $C=12$ and starting index $3$ proves convergence.
The negative first denominator does not enter the tail estimate.

\answer{Part (b) converges to $1$.}

\step{(c) The sequence $c_n=2^{-n}$.}
The expression means the positive reciprocal $1/2^n$.
Fact 2 gives
\[
 \abs{c_n-0}=\frac1{2^n}\leq\frac1{n+1}\leq\frac1n
 \qquad(n\geq1).
\]
The error is thus bounded by $1/n$. Applying Fact 1 proves the
limit directly.

\answer{Part (c) converges to $0$.}

\step{(d) The sequence $t_n=1+2/n$.}
Since $n>0$, every term is defined, and subtracting $1$ gives
\[
 \abs{t_n-1}
 =\abs{1+\frac2n-1}
 =\abs{\frac2n}
 =\frac2n.
\]
For each $\varepsilon>0$, choosing $N>2/\varepsilon$ makes this
error smaller than $\varepsilon$ whenever $n>N$. This is Fact 1
with $C=2$.

\answer{Part (d) converges to $1$.}

\step{(e) The sequence $x_n=73+(-1)^n$.}
If $n$ is even, $(-1)^n=1$ and $x_n=73+1=74$.
If $n$ is odd, $(-1)^n=-1$ and $x_n=73-1=72$.
These values do not merely occur at the start: for any threshold
$N$, choose a positive integer $m>N$. Then $2m>N$ and
$2m+1>N$, and
\[
 x_{2m}=74,\qquad x_{2m+1}=72,\qquad
 \abs{x_{2m}-x_{2m+1}}=\abs{74-72}=2.
\]
Thus Fact 3 applies with $D=2$. No real number can be approached
by every term, because every tail still contains these two values
at distance $2$.

\answer{Part (e) does not converge.}

\step{(f) The sequence $s_n=2^{1/n}$.}
Write $r_n=2^{1/n}$ for the positive real $n$th root of $2$;
thus $r_n>0$ and $r_n^n=2$. We prove the concrete estimate
\[
 0\leq r_n-1\leq\frac1n.
\]
It will imply convergence by Fact 1 without any theorem about
limits of roots or logarithms.

\textbf{First prove the auxiliary Bernoulli inequality.}
For every $x\geq0$ and every positive integer $m$,
\[
 (1+x)^m\geq1+mx.
\]
Fix $x\geq0$. At $m=1$ the two sides are equal. If the inequality
holds at $m$, multiplying by the positive number $1+x$ gives
\begin{align*}
 (1+x)^{m+1}
 &\geq(1+mx)(1+x)\\
 &=1+x+mx+mx^2\\
 &=1+(m+1)x+mx^2\\
 &\geq1+(m+1)x,
\end{align*}
because $mx^2\geq0$. This is the required step at $m+1$,
so induction proves the inequality for the fixed $x$, and hence
for every $x\geq0$.

Apply it with $x=1/n$ and $m=n$. We obtain
\[
 \left(1+\frac1n\right)^n
 \geq1+n\frac1n=2.
\]
For positive real numbers, raising to a positive integer power
preserves strict order. For example, if $u>v>0$, then
$u/v>1$, so the product of $n$ factors $u/v$ is greater than
$1$; multiplication by $v^n>0$ gives $u^n>v^n$.
If $r_n>1+1/n$, this order property would give
\[
 2=r_n^n>\left(1+\frac1n\right)^n\geq2,
\]
a contradiction. Hence $r_n\leq1+1/n$. Also $r_n\geq1$:
if $0<r_n<1$, multiplying $n$ factors each less than $1$
would give $r_n^n<1$, contradicting $r_n^n=2$.

Thus $1\leq r_n\leq1+1/n$, and subtraction of $1$ yields
\[
 \abs{r_n-1}=r_n-1\leq\frac1n.
\]
Fact 1 proves the asserted limit.

\answer{Part (f) converges to $1$.}

\step{(g) The sequence $y_n=n!$.}
For $n\geq1$, the factorial can be written
\[
 n!=n\,(1\cdot2\cdots(n-1)).
\]
At $n=1$, the parenthesized product is empty and equals $1$.
For $n\geq2$, every factor in that product is at least $1$,
so its product is at least $1$. In all cases, therefore,
\[
 n!\geq n.
\]
Now let a real number $B$ and a threshold $N$ be arbitrary.
Choose $n\in\N$ with $n>\max\{N,B,1\}$. Then $n>N$ and
$y_n=n!\geq n>B$. Thus arbitrarily late values exceed every
prescribed real size. Fact 4 rules out convergence to any real
number. This argument establishes the needed nonconvergence
without citing a boundedness theorem.

\answer{Part (g) does not converge to a real limit; its values
grow beyond every real bound.}

\step{(h) The sequence $d_n=(-1)^n n$.}
At an even index $2m$ the value is $2m$, while at an odd index
$2m+1$ the value is $-(2m+1)$. Given any threshold $N$, choose
a positive integer $m>N$. Both indices $2m$ and $2m+1$ are then
greater than $N$, and
\begin{align*}
 \abs{d_{2m}-d_{2m+1}}
 &=\abs{2m-(-(2m+1))}\\
 &=\abs{4m+1}\\
 &=4m+1\geq5>1.
\end{align*}
Here $m\geq1$, so $4m+1$ is positive, justifying the removal
of the absolute value. Every tail has two values separated by
at least $1$. Fact 3 with $D=1$ rules out a real limit.

\answer{Part (h) does not converge. The signs alternate while
the magnitudes increase.}

\step{(i) The sequence $u_n=(-1)^n/n$.}
Although its sign alternates, its distance from $0$ is
\[
 \abs{u_n-0}
 =\frac{\abs{(-1)^n}}{\abs n}
 =\frac1n,
\]
because $\abs{(-1)^n}=1$ and $n>0$. Fact 1 applies with
$C=1$. In contrast with (h), the alternating values now have
small magnitude for every sufficiently large index.

\answer{Part (i) converges to $0$.}

\step{(j) The sequence $u_n=(7n^3+8n)/(2n^3-3)$.}
At $n=1$, the denominator is $2-3=-1$, so it is nonzero.
For $n\geq2$, $n^3\geq8$, giving
$2n^3-3\geq16-3=13>0$. Thus the sequence is defined at every
positive integer.

We test the candidate $7/2$ by exact subtraction:
\begin{align*}
 u_n-\frac72
 &=\frac{2(7n^3+8n)-7(2n^3-3)}{2(2n^3-3)}\\
 &=\frac{14n^3+16n-14n^3+21}{2(2n^3-3)}\\
 &=\frac{16n+21}{2(2n^3-3)}.
\end{align*}
For $n\geq2$, we have $n^3\geq8>3$, so
\[
 2n^3-3\geq n^3>0.
\]
Also $n\geq1$ implies $21\leq21n$ and hence
$16n+21\leq37n$. Using these positive denominator bounds gives
\[
 \abs{u_n-\frac72}
 =\frac{16n+21}{2(2n^3-3)}
 \leq\frac{37n}{2n^3}
 =\frac{37}{2n^2}
 \leq\frac{37}{2n}.
\]
The last step again uses $n^2\geq n>0$. Fact 1, with
$C=37/2$ and starting index $2$, proves convergence.

\answer{Part (j) converges to $7/2$.}

\step{(k) The sequence $u_n=(9n^2-18)/(6n+18)$.}
The denominator $6n+18$ is positive for every $n\geq1$.
We establish a lower estimate that becomes arbitrarily large.
For $n\geq3$, $n^2\geq9$, so $3n^2\geq27>18$ and
\[
 9n^2-18\geq9n^2-3n^2=6n^2>0.
\]
Also $n\geq3$ implies $18\leq6n$, giving
\[
 0<6n+18\leq12n.
\]
First lower the numerator while keeping the positive denominator
fixed, then increase that denominator:
\[
 u_n=\frac{9n^2-18}{6n+18}
 \geq\frac{6n^2}{6n+18}
 \geq\frac{6n^2}{12n}
 =\frac n2.
\]
The second comparison follows from order reversal for positive
reciprocals, followed by multiplication by $6n^2>0$.

Let $B\in\R$ and a threshold $N$ be arbitrary. Choose a positive
integer $n>\max\{N,3,2B\}$. The estimate applies and gives
$u_n\geq n/2>B$, with $n>N$. Fact 4 therefore rules out
convergence to every real candidate.

\answer{Part (k) does not converge to a real limit; its values
grow beyond every real bound.}

\step{(l) The sequence $u_n=\sin(n\pi/2)$.}
Every positive integer has the form $n=4q+r$, with $q\geq0$
an integer and $r\in\{1,2,3,4\}$. Since
$n\pi/2=2q\pi+r\pi/2$, the period $2\pi$ of sine gives
the repeating values
\[
\begin{array}{c|cccc}
 r&1&2&3&4\\ \hline
 \sin(r\pi/2)&1&0&-1&0
\end{array}
\]
The repetition alone suggests nonconvergence; we now check the
arbitrarily late separation required by Fact 3. Given a threshold
$N$, choose $m\in\N$ with $m>N$. Then
$p=4m+1>N$ and $q=4m+3>N$, and periodicity gives
\[
 u_p=\sin(2m\pi+\pi/2)=1,
 \qquad
 u_q=\sin(2m\pi+3\pi/2)=-1.
\]
Their distance is $\abs{1-(-1)}=2$. Fact 3 applies with
$D=2$ and proves nonconvergence.

\answer{Part (l) does not converge.}

\step{(m) The sequence $u_n=\sin(n\pi)$.}
For every integer $n$, the angle $n\pi$ lies at an integer multiple
of a half-turn, so its sine is $0$. Explicitly, if $n=2m$ is
even, then $\sin(n\pi)=\sin(2m\pi)=0$ by periodicity.
If $n=2m+1$ is odd, then
\[
 \sin(n\pi)=\sin(2m\pi+\pi)=\sin\pi=0.
\]
These two cases cover every positive integer. Thus $u_n=0$
for every index, and for every $\varepsilon>0$,
\[
 \abs{u_n-0}=0<\varepsilon
\]
holds at every index. This verifies convergence directly from
the definition.

\answer{Part (m) converges to $0$.}

\step{(n) The sequence $u_n=\sin(2n\pi/3)$.}
Write $n=3q+r$, with $q\geq0$ and $r\in\{1,2,3\}$.
Then $2n\pi/3=2q\pi+2r\pi/3$, giving the period table
\[
\begin{array}{c|ccc}
 r&1&2&3\\ \hline
 \sin(2r\pi/3)&\sqrt3/2&-\sqrt3/2&0
\end{array}
\]
For any threshold $N$, choose a positive integer $m>N$.
The indices $p=3m+1$ and $q=3m+2$ both exceed $N$, and
\[
 u_p=\frac{\sqrt3}{2},\qquad
 u_q=-\frac{\sqrt3}{2}.
\]
Their distance is
\[
 \abs{u_p-u_q}
 =\abs{\frac{\sqrt3}{2}+\frac{\sqrt3}{2}}
 =\sqrt3>0.
\]
Consequently Fact 3 with $D=\sqrt3$ excludes every real limit.

\answer{Part (n) does not converge.}

\step{(o) The sequence $u_n=(\sin n)/n$.}
Here the argument is $n$ radians, with no factor of $\pi$.
We do not need to classify the behavior of $\sin n$ itself.
The elementary trigonometric inequality
$\abs{\sin x}\leq1$ for every real $x$ gives
\[
 \abs{u_n-0}
 =\abs{\frac{\sin n}{n}}
 =\frac{\abs{\sin n}}{n}
 \leq\frac1n,
\]
where division by $n>0$ preserves the inequality.
Fact 1 with $C=1$ proves the limit.

\answer{Part (o) converges to $0$.}

\step{(p) The sequence $u_n=(2^{n+1}+5)/(2^n-7)$.}
The denominator is $2^n-7$. At $n=1$ it
is $2-7=-5$, and at $n=2$ it is $4-7=-3$.
For $n\geq3$, $2^n\geq8$ and therefore $2^n-7\geq1>0$.
These cases show that every term is defined. For example the
first three terms are
\[
 u_1=\frac{4+5}{2-7}=-\frac95,
 \qquad u_2=\frac{8+5}{4-7}=-\frac{13}{3},
 \qquad u_3=\frac{16+5}{8-7}=21.
\]
They do not by themselves determine the limiting behavior.

Subtract the candidate $2$ exactly:
\begin{align*}
 u_n-2
 &=\frac{2^{n+1}+5-2(2^n-7)}{2^n-7}\\
 &=\frac{2\cdot2^n+5-2\cdot2^n+14}{2^n-7}\\
 &=\frac{19}{2^n-7}.
\end{align*}
For $n\geq4$, $2^n\geq16\geq14$, so $7\leq2^n/2$.
It follows that
\[
 2^n-7\geq2^n-\frac{2^n}{2}=\frac{2^n}{2}>0.
\]
Using this bound and Fact 2 yields
\begin{align*}
 \abs{u_n-2}
 &=\frac{19}{2^n-7}\\
 &\leq\frac{19}{2^n/2}
 =\frac{38}{2^n}
 \leq\frac{38}{n+1}
 \leq\frac{38}{n}.
\end{align*}
Fact 1 with $C=38$ and starting index $4$ proves convergence.
It also explains why the negative initial terms and the relatively
large value $u_3=21$ do not change the tail conclusion.

\answer{Part (p) converges to $2$.}

\step{(q) The sequence $a_n=3^n/n!$.}
Both $3^n$ and $n!$ are positive for every $n\in\N$, so
$a_n>0$ and every ratio below is defined. Compare successive
terms using $(n+1)!=(n+1)n!$:
\begin{align*}
 \frac{a_{n+1}}{a_n}
 &=\frac{3^{n+1}}{(n+1)!}\cdot\frac{n!}{3^n}\\
 &=\frac{3\cdot3^n\cdot n!}{(n+1)n!\,3^n}\\
 &=\frac3{n+1}.
\end{align*}
For $n\geq5$, $n+1\geq6$, giving
\[
 0<\frac3{n+1}\leq\frac36=\frac12,
 \qquad
 a_{n+1}\leq\frac12a_n.
\]
We will turn this ratio estimate into an explicit reciprocal bound;
no theorem about ratios of sequences is being invoked.

The value at the starting index is
\[
 C=a_5=\frac{3^5}{5!}
 =\frac{243}{120}=\frac{81}{40}>0.
\]
We claim, for every integer $m\geq0$, that
\[
 a_{5+m}\leq C\,2^{-m}.
\]
At $m=0$, this is equality $a_5=C$. If the estimate holds
at $m$, the ratio bound applies at $n=5+m\geq5$, and gives
\[
 a_{5+(m+1)}=a_{6+m}
 \leq\frac12a_{5+m}
 \leq\frac12C\,2^{-m}
 =C\,2^{-(m+1)}.
\]
This is the next estimate, so induction proves the claim.
Putting $m=n-5$ now yields, for every $n\geq5$,
\[
 0<a_n\leq C\,2^{-(n-5)}.
\]
Fact 2, applied to $m=n-5\geq0$, supplies
\[
 2^{n-5}\geq(n-5)+1=n-4>0.
\]
Taking reciprocals and multiplying by $C>0$ therefore gives
\[
 \abs{a_n-0}=a_n
 \leq C\,2^{-(n-5)}
 \leq\frac{C}{n-4}
 =\frac{81/40}{n-4}.
\]
For any $\varepsilon>0$, choose an integer
$N>\max\{5,4+C/\varepsilon\}$. If $n>N$, then
$n-4>C/\varepsilon$, so the displayed error is less than
$\varepsilon$. This is the shifted estimate of Fact 1, proving
the limit directly.

\answer{Part (q) converges to $0$.}

\step{(r) The sequence $u_n=(1+1/n)^2$.}
Expand by multiplying its two factors, not by assuming a rule
about limits of squares:
\begin{align*}
 u_n
 &=\left(1+\frac1n\right)\left(1+\frac1n\right)\\
 &=1+\frac1n+\frac1n+\frac1{n^2}\\
 &=1+\frac2n+\frac1{n^2}.
\end{align*}
All terms after the leading $1$ are positive. Since
$n^2\geq n>0$, we have $1/n^2\leq1/n$, giving
\[
 \abs{u_n-1}
 =\frac2n+\frac1{n^2}
 \leq\frac2n+\frac1n
 =\frac3n.
\]
Fact 1 with $C=3$ proves convergence.

\answer{Part (r) converges to $1$.}

\step{(s) The sequence $u_n=(4n^2+3)/(3n^2-2)$.}
For all $n\geq1$, $n^2\geq1$, so
$3n^2-2\geq3-2=1>0$. Thus every term is defined and
the denominator has a fixed positive sign.
Subtracting $4/3$ with a common denominator gives
\begin{align*}
 u_n-\frac43
 &=\frac{3(4n^2+3)-4(3n^2-2)}{3(3n^2-2)}\\
 &=\frac{12n^2+9-12n^2+8}{3(3n^2-2)}\\
 &=\frac{17}{3(3n^2-2)}.
\end{align*}
The estimate $n^2\geq1$ also yields
\[
 (3n^2-2)-n^2=2n^2-2=2(n^2-1)\geq0,
\]
so $3n^2-2\geq n^2>0$. Therefore
\[
 \abs{u_n-\frac43}
 =\frac{17}{3(3n^2-2)}
 \leq\frac{17}{3n^2}
 \leq\frac{17}{3n}.
\]
Fact 1 with $C=17/3$ verifies the limit from the definition.

\answer{Part (s) converges to $4/3$.}

\step{(t) The sequence $u_n=(6n+4)/(9n^2+7)$.}
For $n\geq1$, both $6n+4$ and $9n^2+7$ are positive,
so the absolute value of $u_n$ is $u_n$ itself.
We have
\[
 9n^2+7\geq9n^2>0,
 \qquad 6n+4\leq6n+4n=10n,
\]
where $4\leq4n$ uses $n\geq1$. Taking reciprocals in the
denominator comparison and then applying the numerator comparison,
\[
 \abs{u_n-0}
 =\frac{6n+4}{9n^2+7}
 \leq\frac{6n+4}{9n^2}
 \leq\frac{10n}{9n^2}
 =\frac{10}{9n}.
\]
Fact 1 with $C=10/9$ proves convergence to the proposed value $0$.

\answer{Part (t) converges to $0$.}

\subsubsection*{Complete answer list}

In this list, ``does not converge'' means that there is no finite
real limit. Every entry is justified in its corresponding part above.
The growth cases (g) and (k), the persistent oscillations (e), (h),
(l), and (n), and the shrinking oscillations (i) and (o) have been
distinguished by direct estimates and the definition.

\begin{center}
\begin{tabular}{@{}cl@{\qquad}cl@{}}
\toprule
Part&Answer&Part&Answer\\ \midrule
(a)&Converges to $1$&(k)&Does not converge\\
(b)&Converges to $1$&(l)&Does not converge\\
(c)&Converges to $0$&(m)&Converges to $0$\\
(d)&Converges to $1$&(n)&Does not converge\\
(e)&Does not converge&(o)&Converges to $0$\\
(f)&Converges to $1$&(p)&Converges to $2$\\
(g)&Does not converge&(q)&Converges to $0$\\
(h)&Does not converge&(r)&Converges to $1$\\
(i)&Converges to $0$&(s)&Converges to $4/3$\\
(j)&Converges to $7/2$&(t)&Converges to $0$\\
\bottomrule
\end{tabular}
\end{center}

\exercise{7.4}{Rational and irrational terms can have different kinds of limits}
\goal Give \textup{(a)} a sequence of irrational numbers converging to a
rational number, and \textup{(b)} a sequence of rational numbers
converging to an irrational number.

\idea For part \textup{(a)}, divide $\sqrt2$ by larger positive integers.
For part \textup{(b)}, use the lower decimal truncations of $\sqrt2$.
We verify both the nature of every term and convergence.

\step{1. Recall why $\sqrt2$ is irrational}
The irrationality of $\sqrt2$ was established in Section~2.
Here is the short integer argument as well. If
$\sqrt2=p/q$ in lowest terms, with integers $p,q$ and $q>0$,
then $p^2=2q^2$, so $p^2$ is even. The square of an odd integer
$2j+1$ is $4j(j+1)+1$, which is odd. Thus $p$ must be even;
write $p=2j$. Substitution gives
\[
 4j^2=2q^2,\qquad q^2=2j^2.
\]
Now $q^2$ is even, so the same odd-square argument makes $q$ even.
The numerator and denominator would have the common factor $2$,
contrary to lowest terms. Hence $\sqrt2\notin\Q$.

\step{2. Part (a): Irrational terms with rational limit}
Define, for every $n\in\N$,
\[
 x_n=\frac{\sqrt2}{n}.
\]
Since $n\ge1$, division is legitimate. If any term $x_n$ were
rational, multiplication by the rational integer $n$ would make
$nx_n=\sqrt2$ rational. This contradicts Step~1.
Thus every term $x_n$ is irrational.

To prove convergence, let $\varepsilon>0$ and choose an integer
$N>\sqrt2/\varepsilon$ by the Archimedean property. For $n>N$,
\[
 \abs{x_n-0}=\frac{\sqrt2}{n}<\varepsilon.
\]
Therefore $x_n\to0$. The limit is rational because $0=0/1$.

\step{3. Part (b): Define each decimal truncation using a finite maximum}
For a real $t\ge0$, the notation $\lfloor t\rfloor$ denotes the
integer $m$ satisfying
\[
 m\le t<m+1.
\]
To justify existence, choose a positive integer $M>t$ using the
Archimedean property. The finite set
\[
 F=\{j\in\Z:0\le j\le M\text{ and }j\le t\}
\]
is nonempty because it contains $0$. Let $m=\max F$, which exists
for a nonempty finite set of integers. Then $m\le t$.
If $m+1\le t$, the inequalities
$0\le m+1\le t<M$ would give $m+1\in F$, contradicting the
maximality of $m$. Hence $t<m+1$, as required.

This integer is unique. If two integers $m<j$ satisfied their
respective inequalities, then $m+1\le j\le t$ would contradict
$t<m+1$. Interchanging them handles the other possible order.
Thus the floor notation is well-defined here.

Apply this construction to $t=10^n\sqrt2$ and put
\[
 m_n=\left\lfloor10^n\sqrt2\right\rfloor,
 \qquad r_n=\frac{m_n}{10^n}\quad(n\in\N).
\]
Every $m_n$ is an integer and $10^n$ is a positive integer.
Therefore every $r_n$ is rational by the definition of $\Q$.
The defining floor inequalities give
\[
 m_n\le10^n\sqrt2<m_n+1.
\]
Dividing by $10^n>0$ preserves order and yields
\[
 r_n\le\sqrt2<r_n+\frac1{10^n}.
\]
After subtracting $r_n$, this becomes the precise truncation error
estimate
\[
 0\le\sqrt2-r_n<\frac1{10^n}.
\]
In fact the left inequality is strict, because a rational $r_n$
cannot equal the irrational number $\sqrt2$.

\step{4. Prove convergence of the rational truncations}
We may compare the error with $1/n$. The inequality $10^n\ge n$
holds for all positive integers $n$: the base case is $10\ge1$;
if it holds for $n$, then
\[
 10^{n+1}=10\cdot10^n\ge10n\ge n+1,
\]
because $10n-(n+1)=9n-1\ge8>0$ for $n\ge1$.
Induction proves the inequality. Taking reciprocals of positive
quantities gives $1/10^n\le1/n$.

Let $\varepsilon>0$ and choose an integer $N>1/\varepsilon$.
For $n>N$, we obtain
\[
 \abs{r_n-\sqrt2}
 =\sqrt2-r_n<\frac1{10^n}\le\frac1n<\varepsilon.
\]
Thus $r_n\to\sqrt2$, and its limit is irrational by Step~1.

\answer{For \textup{(a)}, take $x_n=\sqrt2/n\to0$, with every
term irrational and the limit rational. For \textup{(b)}, take
$r_n=\lfloor10^n\sqrt2\rfloor/10^n\to\sqrt2$, with every term
rational and the limit irrational.}

\exercise{7.5}{Three differences of large square roots}
\goal Find the limits of
\[
 \text{(a)}\ \sqrt{n^2+1}-n,\qquad
 \text{(b)}\ \sqrt{n^2+n}-n,\qquad
 \text{(c)}\ \sqrt{4n^2+n}-2n.
\]

\idea Multiply each difference by its conjugate divided by that
same conjugate. The difference-of-squares identity removes the
subtraction in the numerator. We then bound the error from the
proposed limit directly, so no later limit laws are needed.

\step{1. Part (a): Expand the product with the conjugate}
Let $a_n=\sqrt{n^2+1}-n$. For $n\ge1$, the square root is
nonnegative and $n>0$, so the conjugate
$\sqrt{n^2+1}+n$ is positive. Dividing by it is valid.
Using $(u-v)(u+v)=u^2-v^2$ gives
\begin{align*}
 a_n
 &=\frac{(\sqrt{n^2+1}-n)(\sqrt{n^2+1}+n)}
         {\sqrt{n^2+1}+n}\\
 &=\frac{(\sqrt{n^2+1})^2-n^2}{\sqrt{n^2+1}+n}\\
 &=\frac{n^2+1-n^2}{\sqrt{n^2+1}+n}\\
 &=\frac1{\sqrt{n^2+1}+n}.
\end{align*}
Since $n^2+1>n^2$ and $n\ge0$, the nonnegative square root
$\sqrt{n^2+1}$ is greater than $n$. Indeed, the contrary inequality
would give $n^2+1\le n^2$ upon squaring nonnegative quantities.
Therefore
\[
 0<a_n=\frac1{\sqrt{n^2+1}+n}\le\frac1{2n}.
\]
Given $\varepsilon>0$, choose an integer $N>1/(2\varepsilon)$.
For $n>N$, this estimate gives
\[
 \abs{a_n-0}=a_n\le\frac1{2n}<\varepsilon.
\]
Thus the limit in part \textup{(a)} is $0$.

\step{2. Part (b): Rationalize and divide by the positive index}
Write $b_n=\sqrt{n^2+n}-n$. Its conjugate is positive for
$n\ge1$, and the same multiplication gives
\begin{align*}
 b_n
 &=\frac{(\sqrt{n^2+n}-n)(\sqrt{n^2+n}+n)}
         {\sqrt{n^2+n}+n}\\
 &=\frac{(n^2+n)-n^2}{\sqrt{n^2+n}+n}\\
 &=\frac n{\sqrt{n^2+n}+n}.
\end{align*}
The radicand factors as $n^2+n=n^2(1+1/n)$. Because $n>0$,
\[
 \sqrt{n^2+n}=n\sqrt{1+\frac1n}.
\]
To check this identity, the right-hand side is nonnegative, and
its square is $n^2(1+1/n)=n^2+n$; uniqueness of the nonnegative
square root then gives the equality. Define
$q_n=\sqrt{1+1/n}$. Substitution and cancellation of $n\ne0$ yield
\[
 b_n=\frac n{nq_n+n}
     =\frac n{n(q_n+1)}
     =\frac1{q_n+1}.
\]
Here $q_n>1$, since $q_n\ge0$ and
$q_n^2=1+1/n>1$; a nonnegative $q_n\le1$ would have square
at most $1$. Hence $q_n+1>2$, so $0<b_n<1/2$.

\step{3. Part (b): Compute the exact error from $1/2$}
All the denominators below are positive. Since $b_n<1/2$,
\begin{align*}
 \abs{b_n-\frac12}
 &=\frac12-\frac1{q_n+1}\\
 &=\frac{q_n+1-2}{2(q_n+1)}\\
 &=\frac{q_n-1}{2(q_n+1)}\\
 &=\frac{(q_n-1)(q_n+1)}{2(q_n+1)^2}\\
 &=\frac{q_n^2-1}{2(q_n+1)^2}\\
 &=\frac{(1+1/n)-1}{2(q_n+1)^2}\\
 &=\frac1{2n(q_n+1)^2}
 \le\frac1{8n}.
\end{align*}
The last inequality uses $(q_n+1)^2\ge2^2=4$ and $n>0$.
For any $\varepsilon>0$, choose an integer $N>1/(8\varepsilon)$.
Then $n>N$ implies
$\abs{b_n-1/2}\le1/(8n)<\varepsilon$.
Therefore the limit in part \textup{(b)} is $1/2$.

\step{4. Part (c): Rationalize and factor the square root}
Put $c_n=\sqrt{4n^2+n}-2n$. Since $n\ge1$, its conjugate
$\sqrt{4n^2+n}+2n$ is positive. Consequently
\begin{align*}
 c_n
 &=\frac{(\sqrt{4n^2+n}-2n)(\sqrt{4n^2+n}+2n)}
         {\sqrt{4n^2+n}+2n}\\
 &=\frac{(4n^2+n)-(2n)^2}{\sqrt{4n^2+n}+2n}\\
 &=\frac{4n^2+n-4n^2}{\sqrt{4n^2+n}+2n}\\
 &=\frac n{\sqrt{4n^2+n}+2n}.
\end{align*}
Now $4n^2+n=n^2(4+1/n)$. The positive-index square-root
identity gives
\[
 \sqrt{4n^2+n}=n\sqrt{4+\frac1n},
\]
because both sides are nonnegative and the square of the right-hand
side is $n^2(4+1/n)=4n^2+n$. Let
$p_n=\sqrt{4+1/n}$. Canceling the nonzero factor $n$ gives
\[
 c_n=\frac n{np_n+2n}
     =\frac n{n(p_n+2)}
     =\frac1{p_n+2}.
\]
Since $p_n\ge0$ and $p_n^2=4+1/n>4$, we have $p_n>2$.
Thus $p_n+2>4$ and $0<c_n<1/4$.

\step{5. Part (c): Compute the exact error from $1/4$}
Using $c_n<1/4$ and then multiplying numerator and denominator
by the positive number $p_n+2$, we obtain
\begin{align*}
 \abs{c_n-\frac14}
 &=\frac14-\frac1{p_n+2}\\
 &=\frac{p_n+2-4}{4(p_n+2)}\\
 &=\frac{p_n-2}{4(p_n+2)}\\
 &=\frac{(p_n-2)(p_n+2)}{4(p_n+2)^2}\\
 &=\frac{p_n^2-4}{4(p_n+2)^2}\\
 &=\frac{(4+1/n)-4}{4(p_n+2)^2}\\
 &=\frac1{4n(p_n+2)^2}
 \le\frac1{64n}.
\end{align*}
Here $(p_n+2)^2\ge4^2=16$, so the denominator is at least
$4n\cdot16=64n$. For any $\varepsilon>0$, choose an integer
$N>1/(64\varepsilon)$. Then $n>N$ gives
$\abs{c_n-1/4}\le1/(64n)<\varepsilon$.
This proves that the limit in part \textup{(c)} is $1/4$.

\answer{The limits are \textup{(a)} $0$, \textup{(b)} $1/2$,
and \textup{(c)} $1/4$. The displayed conjugate identities and
error estimates justify all three values directly.}

